\def\topic{Uniqueness quantifier}
\ifdefined\MyHeaderIncluded
\else
\def\qwertyflag{true}

\ifdefined\MyHeaderIncluded

\else
  \newcommand{\foobar}{FooBar}
  \def\MyHeaderIncluded{true}
% if is till the end of the file

\documentclass[10pt,a4paper]{article}
\usepackage[utf8]{inputenc}
%\usepackage[german]{babel}
%\usepackage[cyrillic]{babel}
\usepackage[T1]{fontenc}
%\usepackage[T2A]{fontenc}
%\usepackage[english, russian]{babel}
\usepackage{amsmath}
\usepackage{amsfonts}
\usepackage{amssymb}
\usepackage{amsthm}
\usepackage{mathtools} %coloneqq

\usepackage{bussproofs} %

%\renewcommand{\thesubsection}{\arabic{subsection}}
%\newtheorem{thm}{Theorem}[section]
%\renewcommand{\thethm}{\arabic{subsection}.\arabic{thm}}

\usepackage{thmtools}
\usepackage{hyperref}
\hypersetup{
    colorlinks=true,
    linkcolor=blue,
    filecolor=magenta,      
    urlcolor=cyan,
    pdftitle={Overleaf Example},
    %pdfpagemode=FullScreen,
    }
%\usepackage{stix-type1}%\turnediota
\theoremstyle{definition}

\newtheorem{constructor}{Constructor}[section]
\newtheorem{iaxiom}{Axiom}
\newtheorem{axiom}{Axiom}
\newtheorem*{axiom*}{Axiom}
\newtheorem{definition}{Definition}[section]
\newtheorem*{definition*}{Definition}
\newtheorem{remark}{Remark}[section]
\newtheorem*{remark*}{Remark}
\newtheorem{recall}{Recall}[section]
\newtheorem*{recall*}{Recall}

\newtheorem{lemma}{Lemma}[section]
\newtheorem{theorem}{Theorem}[section]
\newtheorem*{theorem*}{Theorem}
\newtheorem{note}{Note}[section]
\usepackage{etoolbox}

%\newcommand*{\@Dashv}[2]{%
%  \reflectbox{$\m@th#1#2$}%
%}
\newcommand*{\separation}[3]{% 1)binder, 2)set, 3)condition
  \{#1\in #2: #3\}
}

% https://www.freecodecamp.org/news/simple-regex-tricks-for-beginners-3acb3fa257cb/
%  \\mathcal\{T\} |---> \topology
%  \\mathcal\{([^\}]*)\}   |--->  $1
\newcommand{\AND}{\ \ \&\ \ }
\newcommand{\any}{\mathrm{\ any}}
\newcommand{\borel}{\mathcal{B}}
\newcommand{\can}{\mathrm{can}}
\newcommand{\canonize}{\mathrm{canonize}}
\newcommand{\closure}{\mbox{cl}}
\newcommand{\complex}{\mathbb{C}}
\newcommand{\converges}{\longrightarrow}
\newcommand{\countable}{\mbox{\ countable}}
\newcommand{\ldefined}{\mbox{\ defined}}
\newcommand{\dom}{\mbox{dom}}
\newcommand{\Dom}{\mbox{Dom}}
\newcommand{\Done}{\noindent \textbf{Done.}\\}
\newcommand{\lexists}{\mbox{\ exists}}
\newcommand{\lforall}{\mbox{\ forall}}
\newcommand{\lunique}{\mbox{\ unique}}
\newcommand{\family}{\mathrm{\ family}}
\newcommand{\Fm}{\mathrm{Fm}}
\newcommand{\FuncSymb}{\mathrm{FuncSymb}}
\newcommand{\function}{\mathrm{\ function}}
\newcommand{\functional}{\mbox{\ functional}}
\newcommand{\IFF}{\ \Longleftrightarrow\ }
\renewcommand{\iff}{\ \longleftrightarrow\ }
\newcommand{\IMPLIES}{\ \Longrightarrow\ }
\renewcommand{\implies}{\ \longrightarrow\ }
\newcommand{\Inf}{\mathrm{Inf}}
\newcommand{\injective}{\mbox{\ injective}}
\newcommand{\injection}{\mbox{\ injection}}
\newcommand{\interior}{\mbox{int}}
\newcommand{\intersection}{\bigcap}
\newcommand{\isempty}{\mathrm{\ empty}}
\newcommand{\LB}{\mathrm{LB}}
\newcommand{\nghbset}{\mathcal{N}}
\newcommand{\nonempty}{\mathrm{\ nonempty}}
\newcommand{\Ord}{\mathrm{Ord}}
\newcommand{\power}{\mathcal{P}}
\newcommand{\pr}{\mathrm{pr}}
\newcommand{\fst}{\mathrm{fst}}
\newcommand{\snd}{\mathrm{snd}}
\newcommand{\PredSymb}{\mathrm{PredSymb}}
\newcommand{\Proof}{\noindent \textbf{Proof.}\\}
\newcommand{\ran}{\mathrm{ran}}
\newcommand{\Ran}{\mathrm{Ran}}
\newcommand{\real}{\mathbb{R}}
\newcommand{\relation}{\mathrm{\ relation}}
\newcommand{\set}{\mathrm{\ set}}
\newcommand{\singleton}{\mathrm{\ singleton}}
\newcommand{\semantics}{\mathrm{semantics}}
\newcommand{\semialgebra}{\mathcal{S}}
\newcommand{\Sup}{\mathrm{Sup}}
\newcommand{\surjective}{\mbox{\ surjective\ on\ }}
\newcommand{\Term}{\mathrm{Term}}
\newcommand{\Type}{\mathrm{Type}}
\newcommand{\topology}{\mathcal{T}}
\newcommand{\total}{\mbox{\ total\ on\ }}
\newcommand{\transitive}{\mathrm{\ transitive}}
\newcommand{\union}{\bigcup}
\newcommand{\lTrue}{\mathrm{True}}
\newcommand{\lFalse}{\mathrm{False}}
\newcommand{\lUncertain}{\mathrm{Uncertain}}
\newcommand{\lif}{\mathrm{if}}
\newcommand{\lthen}{\mathrm{then}}
\newcommand{\lelse}{\mathrm{else}}
\newcommand{\lexception}{\mathrm{exception}}
\newcommand{\lundef}{\mathrm{undef}}
\newcommand{\univocal}{\mathrm{\ univocal}}
\newcommand{\univocalat}{\mathrm{\ univocal\ at\ }}
\newcommand{\UB}{\mathrm{UB}}
\newcommand{\Un}{\mathcal{U}}
\newcommand{\Var}{\mathrm{Var}}

\newcommand{\notvarphi}{\phi}
\renewcommand{\phi}{\varphi}

\newcommand{\comment}[1]{\quad\mbox{(#1)}}
% local commands
%\preto\section{\undef\blipo} %etoolbox
\preto\subsection{\undef\blipo} %etoolbox
% https://tex.stackexchange.com/questions/105684/limit-macro-newcommand-scope-to-section
\usepackage{graphicx}
\usepackage{authblk}

\setlength\parindent{0pt} % whole doc of \noindent
\usepackage{soul}
%\usepackage{mathgl}

\fi % guard
\ifdefined\topic
\else
\def\topic{No Titel}
\fi
\title{\topic}
\author[1]{Georgy Dunaev}
\begin{document}
\maketitle
\fi
\section{The unique and its properties}
\subsection{Uniqueness quantifier}
\begin{definition}[uniqueness quantifier]
Uniqueness quantifier for classes:\\
If $Y_1$, $Y_2$ are not free in $\phi$ and they are free for Y in $\phi$, then:
\[
(!Y,\ \varphi)\iff (\forall Y_1\forall Y_2,\ \varphi[Y_1/Y]\land\varphi[Y_2/Y]\implies Y_1=Y_2)
\labeltagd{def}{uniqueness quantifier}
\]
%%(We will not use much this definition.)
\end{definition}
\begin{note}[?]
Definite and indefinite descriptions described here only work properly for sets, not for proper classes. This is inevitable (in classical semantics) due to the nature of descriptors which should return something even if nothing  satisfies the requirements.\\
I have a notion of non--classical models, in which this is not a problem, but then logical rules become much more advanced, and all well--studied field of the first--order logic need to be re--done for such models. Therefore it is too expensive, especially in presence of possibility of using Grothendieck universes for large objects instead of classes. That is why descriptor are designed for sets only: it is enough for potential purposes.\\
\end{note}
\begin{definition}[restricted uniqueness quantifier]
Restricted uniqueness quantifier for sets only
\begin{align*}
(!(y\set),\ \varphi)
&\iff 
(!y,\ (y\set)\land\varphi)
\\
&\iff 
(\forall (y_1\set)\forall (y_2\set),\ \varphi[y_1/y]\land\varphi[y_2/y]\implies y_1=y_2)
\label{def:restricted uniqueness quantifier}
\end{align*}
%(\forall y_1\forall y_2,\ \varphi[y_1/y]\land\varphi[y_2/y]\implies y_1=y_2)
\end{definition}

\subsection{Singleton set}
\begin{definition}[Singleton (functional symbol)]
(side note: because of big amount of side conditions, we state that all different variables are different. 
It should not be possible to use it in such way: "$\{x\}=\cc{x}{x=x}$".
How to determine whether during substitution it can be fixed?
I want it to happen automatically, but i don't want the proof to depend on implementation.
What causes proofs(as list of actions) to be inapplicable in different cases.

Everything is solved by the simple maxima: we don't need developed metatheory, because all those theorems can be proven for a model, and then added to the program if they are useful.

How to implement definitions? see comments in source code here

%Проблема со схемами -- их надо вручную добавлять, код менять.

%Но! Можно задать новый язык формул, с подстановками, требованиями к различности переменных. -- это непонятно: как.

%Боковые условия -- они как дополнительный навесок на теорему.

%(можно, ещё, генерировать, а можно просто рэйзить эксепшены, если что-то не выполняется. Тогда получается язык для определений СХЕМ.

%как логгировать пруф? сохранять состояние последнего успешного шага.
)

)



\[
\{w\} = \cc{x}{x=w}
\]
\end{definition}

\begin{definition}[singleton (a predicate symbol)]
\[
(A\singleton) \iff (\exists!a,\ a\in A)
\]
\end{definition}


\subsection{Properties}
\begin{theorem}[?]
A class constructed with conjugation of a formula and the uniqueness of object satisfying the same formula has at most one element.
\[
(!w,\ w\in \cc{y}{\varphi\land (!y,\ \varphi)})
\clabeltagd{thm}{class with at most one object}
\]
\end{theorem}
\Proof
Assume $w_1\in \cc{y}{\varphi\land (!y,\ \varphi)}$ and $w_2\in \cc{y}{\varphi\land (!y,\ \varphi)}$.\\
To show: $w_1=w_2$.
\[
w_1\in \cc{y}{\varphi\land (!y,\ \varphi)}
\IMPLIES
(w_1\set)
\land
\varphi[w_1/y]
\land
(!y,\ \varphi)
\]
\[
w_2\in \cc{y}{\varphi\land (!y,\ \varphi)}
\IMPLIES
(w_2\set)
\land
\varphi[w_2/y]
\land
(!y,\ \varphi)
\]
\[
(!y,\ \varphi)
\LAND
(w_1\set)\land\varphi[w_1/y]
\LAND
(w_1\set)\land\varphi[w_2/y]
\IMPLIES
(w_1=w_2)
\]
\Done

\begin{theorem}[?]
\[
(A\subseteq B)\land(!x,\ x\in B)\implies (!x,\ x\in A)
\]
\end{theorem}
\Proof
Assume $(A\subseteq B)$ and $(!x,\ x\in B)$.
\[
(!x,\ x\in B)
\IMPLIES
(\forall x_1\forall x_2,\  x_1\in B \land x_2\in B \implies x_1=x_2)
\tag{by \myref{def:uniqueness quantifier}}
\]
The claim: $(\forall x_1\forall x_2,\  x_1\in A \land x_2\in A \implies x_1=x_2)$.\\
Fix $x_1$ and $x_2$.
Assume $x_1\in A$ and $x_2\in A$. 
\[
A\subseteq B
\LAND
x_1\in A
\IMPLIES
x_1\in B
\]
\[
A\subseteq B
\LAND
x_2\in A
\IMPLIES
x_2\in B
\]
\[
(!x,\ x\in B)
\LAND
x_1\in B
\LAND
x_2\in B
\IMPLIES
x_1=x_2
\]
The claim is proven.
\[
(\forall x_1\forall x_2,\ x_1\in A \land x_2\in A \implies x_1=x_2)
\IMPLIES
(!x,\ x\in A)
\tag{by \myref{def:uniqueness quantifier}}
\]
\Done
\ifdefined\qwertyflag
\undef\qwertyflag
\end{document}
\fi